\section{Deleting labels while preserving uniform consequences}\label{sec:transfer}
The closure in Section~\ref{sec:pc-space} reduces a question about
arbitrarily long proofs to one about subspaces and linear maps. The
proof works with these spaces of consequences rather than with invariant certificates. Even an invariant target can belong to the image of an equivariant map without having an invariant preimage in characteristic $p$. Averaging over $S_n$ would therefore miss the issue. The construction below transports the whole image and, for PC, its iterated closure.

A representation, or module, here is a vector space with a linear action
of label permutations. An invariant subspace is preserved by every
such permutation; a quotient by it inherits the action. A
\emph{subquotient} is a quotient of an invariant subspace. The trivial
module is $K$ with every permutation acting as the identity. In
characteristic $p$, the usual averaging operator
$(n!)^{-1}\sum_{g\in S_n}g$ is unavailable once $p\mid n!$.
We therefore construct compatible maps directly, rather than obtain
them by averaging.

There is also a reason not to delete variables by substitution. Setting
all $x_{ij}$ with a deleted row label $i$ to zero turns the totality axiom
$\sum_jx_{ij}-1$ into $-1$. The resulting system is immediately refutable;
it is not the intended smaller instance. Instead we need an operation
which deletes labels on permutation bases while preserving the linear
relations among their subspaces.

The proof separates three requirements: transport the actual coefficient
and multiplication maps, preserve the linear operations forming the
consequence spaces, and reflect target membership. We establish them
below using one deletion operation. A common projector then realises
the reverse direction as an explicit lifting operator. Sections~\ref{sec:block-algebra}--\ref{sec:odd}
construct that projector; the present section explains the contracts
it must satisfy and how they imply the desired transfer.

Here \emph{exactness} means preservation of kernels and images, and hence
of the operations in the PC recurrence. For example, $U+V$ is the image
of $(u,v)\mapsto u+v$, while $U\cap V$ is identified with the kernel of
$(u,v)\mapsto u-v$ on $U\oplus V$. Preserving only dimensions, or only
the spaces generated by the input axioms, would not suffice.

\begin{remark}[Why prime-power blocks are natural]\label{rem:prime-power-counts}
For $q=p^a$ and $0\le j<q$, the images of the integers
$\binom{n+q}{j}$ and $\binom n j$ agree in any field of characteristic
$p$. Indeed, in $K[T]$ repeated use of the characteristic-$p$ binomial
identity gives
\[
 (1+T)^{n+q}=(1+T)^n(1+T^q).
\]
Comparing coefficients below degree $q$ proves the assertion. At
$j=q$ an additional coefficient $1$ appears, so the strict inequality
matters. This explains the persistence of small counting expressions
after a prime-power extension. It is only a counting observation:
ordinary-PC transfer also has to preserve intersections and the
restricted multiplication maps, which require the construction below.
\end{remark}

\subsection{The exact deletion lemma}
For this representation-theoretic construction, $K$ is algebraically
closed of characteristic $p$. Lemma~\ref{lem:field} will return the
conclusions about proof degree to an arbitrary coefficient field.
For a finite set $X$, $K[X]$ denotes the vector space with basis
$e_x$ for $x\in X$. A permutation action on $X$ gives the permutation
module $g e_x=e_{gx}$. The group algebra $K[S_N]$ consists of linear
combinations of label permutations, with composition extended
bilinearly as multiplication. Its elements act by the same expressions
on every representation; an equivariant map commutes with these
expressions. This is the standard module language used below
\cite{AlperinLocal}.

Let $P_b(N)=K[\operatorname{Inj}([b],[N])]$ be the permutation module on injective tuples. Write $\mathcal C_r(N)$ for the finite-dimensional $K[S_N]$-modules annihilated by every group-algebra element which annihilates all $P_b(N)$, $0\le b\le r$. Equivalently, it consists of subquotients of finite direct sums of the $P_b(N)$, $b\le r$. Indeed, their sum is a faithful module for the finite-dimensional quotient algebra; the regular module embeds into a finite sum of copies by evaluation on a basis. This class is closed under submodules, quotients and finite direct sums. It contains the permutation module of every bounded basis system of support at most $r$: each orbit is a quotient of a tuple orbit obtained by choosing an ordering of its support.

Thus $\mathcal C_r(N)$ imposes the operator relations shared by the
tuple spaces up to length $r$. It does not require each vector to
mention at most $r$ labels: a vector in $P_1(N)$ may have nonzero
coefficients at all $N$ labels.

The regular action on the block means that exactly one translation
takes any chosen block label to any other. Every retained label is
fixed by those translations.

Take $q=p^a>r$, write $N=n+q$, and let $A\cong(\Z/p\Z)^a$ act regularly on the last $q$ labels. Its coordinate translations
$\sigma_1,\ldots,\sigma_a$ have order $p$ and fix the retained labels
$[n]$. Choose $c_1,\ldots,c_a\in K$ linearly independent over $\Fp$ and put
\[
 z=\sum_{i=1}^a c_i(\sigma_i-1),\qquad z^p=0,\qquad
 \mathcal D(M)=\ker(z:M\to M)/(\ker z\cap\im z).
\]
Here $1$ is the identity permutation, and multiplication of operators
means composition. The relation $z^p=0$ expresses nilpotence of a
linear operator; it does not assert that $z$ itself is zero.
Permutations of the first $n$ labels commute with $z$, so $\mathcal D(M)$ is an $S_n$-module. Equivariant maps induce the maps of $\mathcal D$.

A nilpotent Jordan chain has basis $v_0,\ldots,v_{\ell-1}$ with
$zv_0=0$ and $zv_j=v_{j-1}$ for $j>0$. Its kernel is $Kv_0$;
for $\ell>1$, this line is already in the image. A length-one chain
has zero image instead. These are statements about the operator $z$;
the chain decomposition need not respect the full label action.

On a Jordan decomposition for the nilpotent operator $z$, this quotient
retains the one-dimensional blocks and discards the longer ones. That
operation is generally not exact. The point of the following theorem is
that the bounded-support representation class, with enough retained labels,
has additional structure making it exact. On permutation bases it then
has the concrete effect of deleting the last block of labels.

All tensor products below are over the working field. Concretely,
$K[X]\otimes Y$ is a list of vectors $y_x\in Y$, written
$\sum_x e_x\otimes y_x$. Labels act on both factors by
$g(e_x\otimes y)=e_{gx}\otimes gy$. A natural identification is one
which commutes with the applicable equivariant maps, rather than merely
matching their dimensions.

\begin{theorem}[Exact deletion]\label{thm:deletion}
For $q>r$ and retained rank $n\ge\max\{0,2r-1\}$ there is a choice of $c_1,\ldots,c_a$, linearly independent over $\mathbb F_p$ in the algebraically closed working field, for which $\mathcal D$ is exact on $\mathcal C_r(n+q)$. It fixes the trivial module. For a support-$u$ basis system $\mathcal U$ with $u<q$, there are natural identifications
\begin{align}
 \mathcal D(K[\mathcal U_{n+q}])&\cong K[\mathcal U_n],\label{eq:permdelete}\\
 \mathcal D(K[\mathcal U_{n+q}]\otimes Y)&\cong
 K[\mathcal U_n]\otimes\mathcal D(Y)\label{eq:tensordelete}
\end{align}
for arbitrary $Y$. For rank-independent invariant coefficient matrices, the first identification transports the actual matrix to its realisation at rank $n$.
\end{theorem}

The algebraic certificate is a single group-algebra element $w$ satisfying
$z^{p-1}wz^{p-1}=z^{p-1}$ on all tuple modules of lengths at most $r$.
Sections~\ref{sec:block-algebra}--\ref{sec:odd} construct the common witness and prove this theorem. The natural matrix and tensor identifications are part of that proof.

The witness is one shared inner-inverse expression; it is not an
inverse to the nilpotent operator on the whole space. Exactness is
asserted for the stated class and retained-rank range, with a suitable
coefficient choice, and does not follow from $z^p=0$ alone.

The tensor identification is required only with a permutation module in its first factor. When applying exactness to a tensor product with an arbitrary subspace, we place that product inside a known bounded-support ambient tensor product. Thus no unstated tensor closure of $\mathcal C_r$ is needed.


\subsection{Transporting the entire consequence space}
Fix a positive degree budget $d$ and write $W_N=P_{N,\le d}$
and $C_d(N)=V_\infty(N)$ for the stabilised PC consequence space.
We first transport the spaces $I_d(N)$ and $C_d(N)$, without choosing
any target. This separates the proof-system calculation from the
question of whether a particular uniform expression belongs to it.

For NS the certificate map has source basis the pairs $(\mu,a)$ with
\[
 \deg\mu+\deg f_{a,N}\le d,
 \qquad (\mu,a)\longmapsto\mu f_{a,N}\in W_N.
\]
Its source has support at most $kd+s$, and $W_N$ has support at most
$kd$. The matrix entry at a target monomial $\nu$ is zero unless
$\mu\mid\nu$, and otherwise is $c(a,\nu/\mu)$. These are the actual
rank-independent invariant entries of the certificate map. The
stability threshold $n_0$ makes each axiom degree independent of rank.
Exactness and matrix naturality therefore identify the embedded image
$\Del(I_d(N))$ with $I_d(n)$.

For PC let $X_N=K[\mathcal V_N]$ and $L_N=P_{N,\le d-1}$.
The multiplication map
\[
       m_N:X_N\otimes L_N\longrightarrow W_N,\qquad
                 x_v\otimes f\longmapsto x_vf
\]
has stable entries. Its source is supported on at most
$k+k(d-1)=kd$ labels. The usable axioms give the image $V_0(N)$
of their coefficient map. The recurrence is
\[
 V_{j+1}(N)=V_j(N)+\im\bigl(m_N|_{X_N\otimes(V_j(N)\cap L_N)}\bigr).
\]
Every stage is a submodule of $W_N$, and every multiplication domain
is a submodule of the fixed support-bounded space $X_N\otimes L_N$.
The intersection is the kernel of $(v,l)\mapsto v-l$ on
$V_j(N)\oplus L_N$. Exactness transports that kernel, the multiplication
image and the sum. The natural tensor identification
\eqref{eq:tensordelete} supplies compatibility with the restricted
multiplication map. Induction thus identifies the embedded
$\Del(V_j(N))$ with $V_j(n)$ for every $j$.

Both ascending chains stabilise because their ambient spaces are finite
dimensional. Choose a stage beyond both stabilisation points. We obtain
\begin{equation}\label{eq:consequence-transport}
 \Del(I_d(N))=I_d(n),\qquad \Del(C_d(N))=C_d(n),
\end{equation}
where equality means equality after the natural identification
$\Del(W_N)\simeq W_n$. This transports every cancellation in a PC
proof, including a low-degree result subsequently multiplied by a
variable. No bound on proof length and no symmetry of individual proof
lines was used.

\subsection{Why deletion also reflects target membership}
Preserving an image is not by itself enough: an exact operation can
annihilate a nonzero quotient. The next lemma supplies the reflection
needed for a general target source.

The detector uses the bilinear coefficient pairing on a permutation
basis, $\langle v,w\rangle=\sum_x v_xw_x$. Its orthogonal space
$S^\perp$ consists of vectors pairing to zero with $S$. This pairing
is nondegenerate over any field; it involves no conjugation or
positivity. A vector represents the functional given by pairing
against it.

\begin{lemma}[Nonzero subquotients survive]\label{lem:detect}
Let $P=K[\mathcal U_N]$ have support at most $u$. Suppose deletion is
exact on $\mathcal C_{2u}(N)$, its block size exceeds $2u$, and the
natural identifications \eqref{eq:permdelete}--\eqref{eq:tensordelete}
hold. If $S\subsetneq S'\subseteq P$ are invariant subspaces, then
$\Del(S'/S)\ne0$.
\end{lemma}
\begin{proof}
The permutation basis gives a nondegenerate dot product on $P$,
equivariant under $S_N$. It induces a pairing
\[
                         S^\perp\otimes(S'/S)\longrightarrow K.
\]
This pairing is surjective: for $v\in S'\setminus S$, a linear
functional vanishing on $S$ can take value one at $v$; the dot product
represents that functional by an element of $S^\perp$.
Let $Y=S'/S$. Tensoring over a field is exact, so the displayed
source lies in $P\otimes Y$, itself a subquotient of $P\otimes P$.
The last permutation module has support at most $2u$, justifying
every use of exactness. If $\Del(Y)=0$, permutation-first-factor tensor
compatibility gives $\Del(P\otimes Y)=0$. Exactness then takes the
surjective pairing to a surjection from zero onto
$\Del(K)\simeq K\ne0$, a contradiction.

We use this surjective pairing, not a trace or an averaged inner
product. In characteristic $p$, composing a pairing with coevaluation
can yield a vanishing dimension scalar and would not prove the lemma.
\end{proof}

The invariant-target case has a simpler detector: the trivial module
$K$ is one-dimensional, so a nonzero map from it has a one-dimensional
image. It has no proper nonzero submodule, which is the meaning of
\emph{simple} in the following argument.

\begin{lemma}[An invariant target needs no doubled budget]\label{lem:constant}
Let $j:K\to W$ and $H\subseteq W$ be equivariant and belong to the
exactness class. Then $j(1)\in H$ if and only if its deleted image
belongs to $\Del(H)\subseteq\Del(W)$.
\end{lemma}
\begin{proof}
Membership is the vanishing of $K\xrightarrow{j}W\to W/H$.
Its image is either zero or isomorphic to the simple trivial module
$K$. Exactness preserves this image, and $\Del(K)\simeq K\ne0$.
Thus the composite vanishes exactly when its deleted composite does.
The vector $j(1)$ may be any invariant target, not just the constant
polynomial one.
\end{proof}

\subsection{The actual uniform target map}\label{sec:target-transfer}
For indexed targets, let $\mathcal B$ be the fixed index basis and
let $\gamma(b,\mu)$ be its stable coefficient rule. Write $\delta$ for the largest stabilised target degree, taking
$\delta=0$ for an empty request. At a rank where
all target degrees are at most $d$, define
\[
 J_N:K[\mathcal B_N]\longrightarrow W_N,\qquad
                [b]\longmapsto G_b^{(N)}.
\]
Its entry in row $\mu$, column $b$ is precisely $\gamma(b,\mu)$.
Equivariance and rank independence are the hypotheses of the natural
matrix identification in Theorem~\ref{thm:deletion}. Consequently
$\Del(J_N)$ is the actual smaller target map $J_n$.
This statement concerns the prescribed full polynomials $G_b^{(n)}$;
we do not replace a target by its degree-$d$ truncation.

Let $H_N$ be either $I_d(N)$ or $C_d(N)$. Every target belongs to
$H_N$ exactly when $\im J_N\subseteq H_N$. Put
\[
          S_N=H_N+\im J_N,\qquad Y_N=S_N/H_N.
\]
The condition is $Y_N=0$. Exactness, target-map naturality and
\eqref{eq:consequence-transport} identify its deletion with
\[
          Y_n=(H_n+\im J_n)/H_n.
\]
The polynomial permutation basis of $W_N$ has support at most $kd$.
Lemma~\ref{lem:detect}, applied with this ORIGINAL support bound,
reflects vanishing of $Y_N$ once the budget covers $2kd$.
The affine budget
\[
                            R(d)=2kd+s+t
\]
also covers the certificate source ($kd+s$), the target source ($t$)
and every multiplication domain. It therefore proves membership in
both directions. The target source need not be simple; reflection
comes from a subquotient of the TARGET ambient polynomial space.
Doubling $R(d)$ once again would be unnecessary: the pairing uses
$W_N\otimes W_N$, whose support is $2kd$.

\begin{proof}[Proof of Theorem~\ref{thm:periodicity}]
First extend scalars to an algebraic closure, as permitted by
Lemma~\ref{lem:field}. The degree of every indexed target stabilises
beyond $n_*=\max\{n_0,1,kD_G+t\}$. If $d$ is smaller than the
largest stable target degree, the all-target predicate is false at both
ranks. Otherwise the full map $J$ just described applies.
Take $N=n+q$, $q=p^a>R(d)$ and retained rank
$n\ge\max\{n_*,2R(d)-1\}$. Exact deletion transports $H_N$, $J_N$
and their sum and quotient. The detection lemma yields
\[
 \bigl(\forall b\in\mathcal B_N,\ G_b^{(N)}\in H_N\bigr)
 \quad\Longleftrightarrow\quad
 \bigl(\forall b\in\mathcal B_n,\ G_b^{(n)}\in H_n\bigr).
\]
Field descent holds for each actual target separately, so the equivalence
holds over the original field. Repeated steps connect any eligible
ranks in a common residue class and prove the stated periodicity.
For one fixed global target orbit, replace $\mathcal B$ by that orbit;
at eligible ranks, equivariance makes every representative equivalent.

At degree zero the stabilised targets are either all zero, or are
scalars with at least one nonzero value. In the latter case every target
is derivable in degree zero exactly when the original axioms include a
nonzero constant. If a target has positive stable degree, degree-zero
all-target membership fails. These are fixed predicates beyond $n_*$
and give the separate zero-level assertion.

For a fully invariant target use $j_N:K\to W_N$, $j_N(1)=G^{(N)}$,
in place of $J_N$. Lemma~\ref{lem:constant} reflects membership without
the tensor pairing. The budget reduces to $kd+s$, with stable onset
$\max\{n_0,1,kD_G\}$. Taking $G^{(N)}=1$ gives the sharper refutation
special case. No target has been added to the axiom system.
\end{proof}

\subsection{A lifting operator for bounded-degree consequences}
Fix $d\ge1$, write $W_n=P_{n,\le d}$ and use a block
$q=p^a>r=kd+s$ in the deletion range. Let
$\iota:W_n\to W_N$ include retained monomials and
$\rho:W_N\to W_n$ discard all others, where $N=n+q$.
We seek a common projector lift satisfying
\begin{equation}\label{eq:newcontract}
 \Lift:W_n\longrightarrow W_N,\qquad
 \deg\Lift(f)\le\deg f,\qquad\Lift(1)=1,\qquad
 \Lift(V_j(n))\subseteq V_j(N)\quad\text{for every }j.
\end{equation}
It will also satisfy $\rho\Lift=\mathrm{id}$. The previously proved
transfer gives
\begin{equation}\label{eq:newiff}
 1\in C_d(n)\quad\Longleftrightarrow\quad1\in C_d(N).
\end{equation}
The operator below realises the upward direction while preserving the
whole PC closure. Indexed targets still require the target-map
argument above; their projector lifts need not equal the prescribed
larger targets individually.

\subsection{The lift from a common projector}\label{lift-sec:formula}
We now give the map, leaving the construction of its certificate to the
next sections. Temporarily allow a coefficient extension $E/K$ and use
the same notation for the extended polynomial spaces. Section~\ref{lift-sec:descent}
returns the lifting statement to $K$.

Identify the added block $B$ with $A=\Fp^a$, acting on it by translations
and fixing $[n]$. Let $\sigma_1,\ldots,\sigma_a$ be the coordinate
translations. Choose $c_1,\ldots,c_a\in E$ linearly independent over
$\Fp$ and set
\[
 z=\sum_{i=1}^a c_i(\sigma_i-1)\in E[S_N],\qquad z^p=0.
\]
The group algebra $E[S_N]$ consists of linear combinations of label
permutations; it acts on polynomials by the corresponding linear
combinations of variable renamings. The permutations commute on the
block, and $\sigma_i^p=1$, which proves $z^p=0$.

The finite algebraic target is one element $w\in E[S_N]$ for which
\begin{equation}\label{lift-eq:witness}
 u w u=u\quad\text{on every }P_b(N)=E[\Inj([b],[N])],\quad b\le r,
 \qquad u=z^{p-1}.
\end{equation}
Thus the same expression must work on every tuple length up to the
support budget. This inner-inverse identity does not make the
nilpotent operator $u$ invertible. For characteristic two put
\[
                         e=(1-zw)(1-wz).
\]
For general $p$, put $m=p-1$ and
\begin{equation}\label{lift-eq:projector}
 S_0=1,\qquad S_{j+1}=(1-z^jwz^{m-j})S_j\ (0\le j<m),
 \qquad e=(1-z^mw)S_m.
\end{equation}
An idempotent, or projector, is an operator satisfying $e^2=e$.
For $v\in\ker z$, the constructed $e$ will supply the representative
$ev$ of its deletion class. Its idempotence and this identification
are proved below from the common witness and the block action.

The two formulas agree at $p=2$. The general formula requires a further
property of the block action: every nilpotent Jordan chain on a tuple
module has length one or $p$. Nilpotence $z^p=0$ alone would also allow
intermediate lengths, for which this projector need not work. The
block decomposition establishes this chain property; together with
the common witness it gives the projector identities in
Section~\ref{sec:odd}. Once these inputs are
proved, the formulas are explicit finite algebra expressions. The lifting operator is
\begin{equation}\label{lift-eq:actual-lift}
 \boxed{\quad\Lift_{n,N,d}(f)=e\,\iota(f).\quad}
\end{equation}
Equivalently, writing $e=\sum_{g\in S_N}a_g g$,
\[
                    \Lift(f)=\sum_{g\in S_N}a_g\,g\bigl(\iota(f)\bigr).
\]
This explains why degree does not increase: inclusion and every
renaming preserve degree. The difficult point is to choose the
\emph{common} coefficients $a_g$ so that all consequence spaces are
preserved. An arbitrary sum of renamings does not have that property.

\subsection{Compatibility with the entire PC closure}\label{lift-sec:verify}
The common projector constructed in Sections~\ref{sec:block-algebra}--\ref{sec:odd}
satisfies $e^2=e$, $ze=ez=0$ and, on $\ker z$,
\begin{equation}\label{lift-eq:onker}
                       ev=v-z^{p-1}wv.
\end{equation}
It identifies $\Del(M)$ with $eM$ naturally as vector spaces.
Together with the retained-coordinate description of deletion,
\eqref{eq:consequence-transport} and its stagewise proof give
\begin{equation}\label{lift-eq:stage-transport}
 \rho\bigl(V_j(N)\cap\ker z_{W_N}\bigr)=V_j(n)\quad\text{for every }j.
\end{equation}
Thus the exactness and multiplication-map work need not be repeated.

\begin{proposition}[The lifting formula preserves and reflects consequences]\label{lift-prop:lift}
Assume the two facts in Section~\ref{lift-sec:formula}. Then
$\Lift=e\iota$ preserves degree, $\rho\Lift=\mathrm{id}$,
$\Lift(1)=1$, and
\[
                  \Lift(f)\in V_j(N)\quad\Longleftrightarrow\quad
                  f\in V_j(n).
\]
The same equivalence holds for the stable spaces $C_d$.
\end{proposition}
\begin{proof}
Every retained monomial is fixed by the block action, so $z\iota(f)=0$.
Equation~\eqref{lift-eq:onker} gives
\[
                  \Lift(f)=\iota(f)-z^{p-1}w\iota(f).
\]
The image of $z$ has no retained component: $z$ kills retained basis
vectors and preserves the span of moved ones. Hence $\rho\Lift(f)=f$.
On constants, every permutation acts trivially, so the correction term
vanishes and $\Lift(1)=1$.

If $f\in V_j(n)$, equation~\eqref{lift-eq:stage-transport} supplies
$v\in V_j(N)\cap\ker z$ with $\rho v=f$. The permutation-space
identification says that $v-\iota(f)\in\ker z\cap\im z$. Since
$e\im z=0$, we have $\Lift(f)=ev$. The space $V_j(N)$ is invariant
under all label permutations, hence under $e$, and $ev\in V_j(N)$.
Conversely, $\Lift(f)\in V_j(N)$ is killed by $z$, so
\eqref{lift-eq:stage-transport} and $\rho\Lift(f)=f$ give $f\in V_j(n)$.
Degree preservation follows from the renaming formula. Pass to the
unions for the stable statement.
\end{proof}

This proves both the forward lifting implication and
\eqref{eq:newiff}. The image of the lift need not be all of $C_d(N)$:
the larger universe has many additional consequences. Every smaller consequence has this projector lift. In the transfer range, for $n\ge n_*$ and a fully
invariant target $G$ whose stabilised degree is at most $d$, naturality of its map from $K$ gives
$\Lift(G^{(n)})=G^{(N)}$: the target map commutes with the common projector,
which acts as the identity on $K$. More concretely, invariance gives
$zG^{(N)}=0$ and $eG^{(N)}=G^{(N)}$, while stable coefficients give
$\rho G^{(N)}=G^{(n)}$. The retained-vector identification and the argument
in Proposition~\ref{lift-prop:lift} therefore yield $e\iota(G^{(n)})=G^{(N)}$. Hence invariant targets, including
$1$, are exactly preserved. For indexed targets this equality need
not hold. The image-subquotient argument above proves transfer of
prescribed indexed targets without assuming it.

The common action also explains how individual inference steps are
realised. For a retained axiom index $a$, naturality of its coefficient
map $F_N$ gives
\begin{equation}\label{lift-eq:axiomlift}
                 \Lift(f_{a,n})=F_N(e_A\iota_A(a)).
\end{equation}
For $f\in V_j(n)\cap P_{n,\le d-1}$, act on the joint tensor input:
\begin{equation}\label{lift-eq:jointlift}
 e_{X_N\otimes L_N}(\iota(x_v)\otimes\iota(f))
       =\sum_y x_y\otimes g_y,
 \qquad X_N=K[\mathcal V_N],\quad L_N=P_{N,\le d-1}.
\end{equation}
The tensor identification in \eqref{eq:tensordelete} places each
$g_y$ in $V_j(N)\cap L_N$. Equivariance of multiplication then gives
\begin{equation}\label{lift-eq:liftmul}
                         \Lift(x_vf)=\sum_y x_yg_y.
\end{equation}
The same permutation acts on both tensor factors; replacing this joint
action by $e_X\otimes e_L$ would lose that compatibility. Storing
closure-basis derivations and expressing the $g_y$ in those bases
recovers a literal PC proof, with no proof-length claim.

\subsection{A one-label lifting example}\label{lift-sec:small}
Work in characteristic two. Retain the label $a$ and add $b,c$.
Let $z=1+(bc)$ and $w=(ab)$ act on the linear variables
$x_a,x_b,x_c$. On their three-dimensional span,
\[
 z=\begin{pmatrix}0&0&0\\0&1&1\\0&1&1\end{pmatrix},\qquad
 w=\begin{pmatrix}0&1&0\\1&0&0\\0&0&1\end{pmatrix}.
\]
Direct multiplication gives $zwz=z$ and
\[
                       e=(1-zw)(1-wz)
                         =\begin{pmatrix}1&1&1\\1&1&1\\1&1&1\end{pmatrix}.
\]
Consequently
\begin{equation}\label{lift-eq:smalllift}
                  \Lift(\alpha+\beta x_a)
                      =\alpha+\beta(x_a+x_b+x_c).
\end{equation}
Its retained coordinate is the original polynomial, and $\Lift(1)=1$.

Take the uniform degree-one axiom $\sum_{i<m}x_i=0$ at every universe size $m$.
At size one the axiom is $x_a$; at size three it is $x_a+x_b+x_c$.
Formula~\eqref{lift-eq:smalllift} sends the former exactly to the latter and
lifts every degree-one consequence. Plain inclusion would leave $x_a$,
which is not a consequence of the larger system: the assignment
$(x_a,x_b,x_c)=(1,1,0)$ satisfies its sum but not $x_a=0$.

Here the invariant target $G^{(m)}=\sum_{i<m}x_i$ is transported to
the prescribed larger target. In contrast, the indexed expression
$G_a^{(m)}=x_a$ has projector lift $x_a+x_b+x_c$, not $x_a$. The map
is linear but is not a ring homomorphism and does not preserve every
indexed expression literally. This calculation illustrates the distinction but lies outside the
indexed transfer range: for $k=d=t=1$ and $s=0$, that theorem uses
$R(d)=3$, whereas here $q=2$ and $n=1$. It is therefore no
counterexample to indexed periodicity.

\subsection{Descent to the coefficient field}\label{lift-sec:descent}
Lemma~\ref{lem:field} returns every target-membership equivalence to
the coefficient field $K$. A concrete forward lift can also be descended:
all its coefficients lie in a finite extension $E/K$. Choose a $K$-linear
map $\lambda:E\to K$ with $\lambda(1)=1$ and apply it coefficientwise:
\begin{equation}\label{lift-eq:fieldlift}
                  \widehat\Lift(f)=\lambda_*(\Lift_E(f)).
\end{equation}
Since $V_j(N;E)=E\otimes_K V_j(N;K)$, this preserves membership in each
closure stage, fixes $1$, preserves degree and satisfies
$\rho\widehat\Lift=\mathrm{id}$. Recover a proof from the $K$-closure
bases; coefficientwise projection of individual proof steps need not
preserve their scalar inferences. Reflection of prescribed target
membership follows from Lemma~\ref{lem:field}, not from a claim that
$\widehat\Lift$ reflects every individual consequence.
